\documentclass[11pt]{article}

% Change "review" to "final" to generate the final (sometimes called camera-ready) version.
% Change to "preprint" to generate a non-anonymous version with page numbers.
\usepackage[review]{acl}

% Standard package includes
\usepackage{times}
\usepackage{latexsym}

% For proper rendering and hyphenation of words containing Latin characters (including in bib files)
\usepackage[T1]{fontenc}
% For Vietnamese characters
% \usepackage[T5]{fontenc}
% See https://www.latex-project.org/help/documentation/encguide.pdf for other character sets

% This assumes your files are encoded as UTF8
\usepackage[utf8]{inputenc}

% This is not strictly necessary, and may be commented out,
% but it will improve the layout of the manuscript,
% and will typically save some space.
\usepackage{microtype}

% This is also not strictly necessary, and may be commented out.
% However, it will improve the aesthetics of text in
% the typewriter font.
\usepackage{inconsolata}

%Including images in your LaTeX document requires adding
%additional package(s)
\usepackage{graphicx}

% If the title and author information does not fit in the area allocated, uncomment the following
%
%\setlength\titlebox{<dim>}
%
% and set <dim> to something 5cm or larger.

\title{Mapudungun MT}

% Author information can be set in various styles:
% For several authors from the same institution:
% \author{Author 1 \and ... \and Author n \\
%         Address line \\ ... \\ Address line}
% if the names do not fit well on one line use
%         Author 1 \\ {\bf Author 2} \\ ... \\ {\bf Author n} \\
% For authors from different institutions:
% \author{Author 1 \\ Address line \\  ... \\ Address line
%         \And  ... \And
%         Author n \\ Address line \\ ... \\ Address line}
% To start a separate ``row'' of authors use \AND, as in
% \author{Author 1 \\ Address line \\  ... \\ Address line
%         \AND
%         Author 2 \\ Address line \\ ... \\ Address line \And
%         Author 3 \\ Address line \\ ... \\ Address line}

\author{First Author \\
  Affiliation / Address line 1 \\
  Affiliation / Address line 2 \\
  Affiliation / Address line 3 \\
  \texttt{email@domain} \\\And
  Second Author \\
  Affiliation / Address line 1 \\
  Affiliation / Address line 2 \\
  Affiliation / Address line 3 \\
  \texttt{email@domain} \\}

%\author{
%  \textbf{First Author\textsuperscript{1}},
%  \textbf{Second Author\textsuperscript{1,2}},
%  \textbf{Third T. Author\textsuperscript{1}},
%  \textbf{Fourth Author\textsuperscript{1}},
%\\
%  \textbf{Fifth Author\textsuperscript{1,2}},
%  \textbf{Sixth Author\textsuperscript{1}},
%  \textbf{Seventh Author\textsuperscript{1}},
%  \textbf{Eighth Author \textsuperscript{1,2,3,4}},
%\\
%  \textbf{Ninth Author\textsuperscript{1}},
%  \textbf{Tenth Author\textsuperscript{1}},
%  \textbf{Eleventh E. Author\textsuperscript{1,2,3,4,5}},
%  \textbf{Twelfth Author\textsuperscript{1}},
%\\
%  \textbf{Thirteenth Author\textsuperscript{3}},
%  \textbf{Fourteenth F. Author\textsuperscript{2,4}},
%  \textbf{Fifteenth Author\textsuperscript{1}},
%  \textbf{Sixteenth Author\textsuperscript{1}},
%\\
%  \textbf{Seventeenth S. Author\textsuperscript{4,5}},
%  \textbf{Eighteenth Author\textsuperscript{3,4}},
%  \textbf{Nineteenth N. Author\textsuperscript{2,5}},
%  \textbf{Twentieth Author\textsuperscript{1}}
%\\
%\\
%  \textsuperscript{1}Affiliation 1,
%  \textsuperscript{2}Affiliation 2,
%  \textsuperscript{3}Affiliation 3,
%  \textsuperscript{4}Affiliation 4,
%  \textsuperscript{5}Affiliation 5
%\\
%  \small{
%    \textbf{Correspondence:} \href{mailto:email@domain}{email@domain}
%  }
%}

\newcommand{\todo}[1]{\textbf{[TODO: #1]}}

\begin{document}
\maketitle
\begin{abstract}
Massive multilingual translation models have advanced language technologies, yet many languages have still been left out. One such case is Mapudungun, an indigenous language spoken by 200,000+ Mapuche people in Chile and Argentina. Prior work on this language for machine translation has been conducted minimally before the advent of modern large language models. Mapudungun's grammar is polysynthetic and agglutinative, which has been shown to cause issues for standard tokenization techniques. In this work, we fine-tune three sizes of Meta's NLLB-200 with a comparison of 8 tokenization methods, including our novel method Morfessor-VC, which takes  the automatic morpheme segmenter Morfessor and constrains its vocabulary to NLLB's vocabulary. We found NLLB's smallest model with 600M parameters using our method is on par with NLLB's largest 3.3B parameter model with standard BPE, and beats modern LLMs by nearly 30 chrF++. This paper provides a technical foundation that would enable the creation of high-quality language tools for pedagogy, social networks, and language immersion for the Mapuche community.
\end{abstract}

\section{Introduction}

Mapudungun (also \textit{Mapuzugun} or \textit{Chedungun}) is spoken by more than 200,000 Mapuche people across southern Chile and Argentina \cite{duan2020avenue}, making it one of the largest indigenous languages in South America by speaker count. Despite this, Mapudungun remains absent from nearly all commercial and open-source machine translation (MT) systems. The language's polysynthetic, agglutinative morphology—where a single verb form can encode tense, aspect, subject, object, evidentiality, and spatial deixis—poses well-documented challenges for standard subword tokenization \cite{rust2021good,mielke2021between,petrov2023language}. Byte-pair encoding (BPE) \cite{sennrich2016bpe}, the dominant tokenization strategy in multilingual models, is not designed to respect morpheme boundaries, and high fertility rates on agglutinative languages have been shown to correlate with degraded translation quality \cite{ahia2023all,banerjee2018morpheme}.

Prior work on Mapudungun NLP is sparse. \citet{duan2020avenue} introduced the AVENUE corpus and established early sequence-to-sequence baselines. \citet{pendas2023} and \citet{liranmt2025} have since explored neural approaches, and \citet{chandia2022} studied automatic morphological analysis. None of this prior work has fine-tuned large-scale multilingual models on Mapudungun, nor systematically studied how tokenization choices interact with model capacity for this language.

We address both gaps. Our contributions are:

\begin{itemize}
  \item The first fine-tuning study of Meta's NLLB-200 \cite{nllbteam2022} on Mapudungun--Spanish translation, across three model sizes (600M, 1.3B, 3.3B parameters) and both translation directions.

  \item A systematic ablation of eight tokenization strategies, ranging from standard NLLB BPE to Morfessor-based segmentation \cite{virpioja2013morfessor} and SentencePiece UnigramLM \cite{kudo2018subword}.

  \item \textbf{Morfessor-VC}, a novel tokenization method that runs Morfessor segmentation and then constrains the resulting vocabulary to tokens already present in NLLB's pretrained vocabulary, preserving embedding alignment while improving morpheme boundary coverage.

  \item A comprehensive evaluation showing that Morfessor-VC with the 600M model achieves 43.2 chrF++ on Mapudungun$\rightarrow$Spanish, matching our standard BPE 3.3B baseline (42.9 chrF++), and outperforming frontier LLMs (Aya Expanse 8B: 15.9 chrF++) by over 27 points.
\end{itemize}

These results demonstrate that morphology-aware tokenization can substitute for a 5$\times$ increase in model parameters for Mapudungun MT, with direct implications for low-resource, compute-constrained deployment in Mapuche community settings. An additional complication is that the AVENUE corpus contains substantial code-switching: approximately 24\% of Mapudungun utterances contain Spanish words, reflecting natural bilingual speech in the Mapuche community. We analyze how this affects translation quality in both directions in Section~\ref{sec:cs}.

% TODO: add code-switching sentence to Introduction (see below)

\section{Related Work}
% TODO: Mapudungun NLP: Duan et al. (2020), Pendas et al. (2023), Ahumada et al. (2022), Chandía (2022), Lira et al. (2025)
% TODO: Morphology-aware tokenization: Morfessor (Virpioja 2013), UnigramLM (Kudo 2018), Banerjee & Bhattacharyya (2018), Gowda & May (2020), Ataman & Federico (2018)
% TODO: Low-resource MT with multilingual models: NLLB (Costa-jussà 2022), AmericasNLP shared tasks

\todo{Draft Related Work}

\section{Data}

We use the AVENUE corpus \cite{duan2020avenue}, a Mapudungun--Spanish parallel corpus collected through oral history interviews with native Mapuche speakers in Chile, transcribed and aligned using the ELAN annotation tool. The corpus contains bilingual conversational speech covering domains such as traditional medicine, cultural practices, and personal narrative.

\paragraph{Extraction.} The raw AVENUE data is structured as parallel ELAN blocks. We extract sentence-level pairs using the \textit{blocks} segmentation approach, which aligns ELAN annotation blocks directly. After extracting parallel pairs, we apply a cleaning pipeline that removes empty lines, mismatched lengths, and degenerate segments. The resulting corpus is split into train / dev / test sets of 55,452 / 1,581 / 9,382 sentence pairs, respectively.

\paragraph{Code-switching.} A notable property of the AVENUE corpus is widespread code-switching: approximately 24\% of Mapudungun utterances contain embedded Spanish words or phrases, as marked by ELAN \textsc{<SPA>} tags in the original transcriptions. This reflects natural bilingual speech patterns in contemporary Mapuche communities rather than transcription artifacts. Our cleaning pipeline retains these mixed-language utterances; we analyze their effect on translation quality in Section~\ref{sec:cs}.

\paragraph{Data statistics.} The Mapudungun training side contains 656,969 whitespace-delimited tokens (avg.\ 11.8 tokens/sentence); the Spanish side contains 888,476 tokens (avg.\ 16.0 tokens/sentence). The higher Spanish token count reflects that Mapudungun's polysynthetic morphology encodes in a single word what Spanish expresses across multiple words.

\todo{Insert data statistics table (token counts, type counts, avg.\ sentence length per split)}

\section{Tokenization Methods}
\label{sec:tokenization}

A core challenge in fine-tuning NLLB-200 on Mapudungun is the mismatch between NLLB's pretrained BPE vocabulary—optimized for 200 languages with Spanish heavily represented—and Mapudungun's polysynthetic morphology. Naive fine-tuning feeds the model unsegmented Mapudungun text, which NLLB's internal tokenizer fragments unpredictably. We compare eight tokenization strategies that differ in how they pre-segment Mapudungun before fine-tuning. All conditions use NLLB's standard tokenizer for the Spanish side; only the Mapudungun side is varied. Pre-segmented text uses the \texttt{@@} boundary marker convention (e.g., \textit{mapu\@@\@ dun\@@\@ gun}) so that desegmentation is applied at inference time before scoring.

\subsection{Standard BPE (NLLB)}

The baseline condition feeds raw, unsegmented Mapudungun text directly to NLLB's built-in SentencePiece tokenizer without any pre-processing. This is the standard fine-tuning setup for NLLB and serves as our primary comparison point.

\subsection{Joint BPE}

Following \citet{duan2020avenue}, we train a shared SentencePiece BPE model \cite{kudo2018sentencepiece} on the concatenated Mapudungun and Spanish training data with a vocabulary size of 5,000. Joint training encourages shared subword representations across languages, which can aid cross-lingual transfer but may produce poor segmentations for the morphologically richer language.

\subsection{Mono BPE}

We train a Mapudungun-only SentencePiece BPE model, setting the vocabulary size via the 95\%-character-coverage heuristic of \citet{gowda2020finding}: we find the number of character types needed to cover 95\% of character occurrences, then multiply by 10 to estimate the subword vocabulary size. This yields a vocabulary focused on Mapudungun morphology without interference from Spanish.

\subsection{Optuna BPE}

To address the objection that our results might be sensitive to vocabulary size selection, we use Optuna \cite{akiba2019optuna} to tune the BPE vocabulary size over 15 trials, optimizing chrF++ on the development set using a 600M NLLB model fine-tuned for 3 epochs per trial. The optimal vocabulary size found was 10,954. Notably, chrF++ was largely flat across a wide range of vocabulary sizes (1,431--14,517 all achieved 45.72 chrF++), suggesting BPE vocabulary size is not a critical hyperparameter for this language pair.

\subsection{Morfessor}

We apply Morfessor 2.0 \cite{virpioja2013morfessor}, an unsupervised morpheme segmenter based on minimum description length, to the Mapudungun training corpus. Morfessor learns morpheme boundaries without linguistic supervision, producing segmentations that reflect statistical regularities in the character sequences. Boundaries are marked with \texttt{@@}. The Spanish side is left unsegmented.

\subsection{Morfessor-VC (Our Method)}
\label{sec:morfessor-vc}

Standard Morfessor segmentation introduces boundaries that do not align with NLLB's pretrained vocabulary, causing \textit{double tokenization}: NLLB's internal tokenizer re-splits the already-segmented morphemes, producing fragmented representations that were not seen during pretraining.

\textbf{Morfessor-VC} (Vocabulary-Constrained) is our proposed solution. Starting from Morfessor segmentation, we greedily merge adjacent morphemes whose concatenation corresponds to a single token in NLLB's SentencePiece vocabulary. The algorithm processes each word left-to-right:

\begin{enumerate}
  \item Given morphemes $[m_0, m_1, \ldots, m_k]$ within a word, consider the candidate merge $m_i \oplus m_{i+1}$.
  \item If $\text{▁}(m_i \oplus m_{i+1})$ or $(m_i \oplus m_{i+1})$ is a single piece in NLLB's vocabulary, merge and repeat from $m_i$.
  \item Otherwise, retain the boundary and advance to $m_{i+1}$.
\end{enumerate}

The result preserves morpheme boundaries only where the split is both linguistically motivated (by Morfessor) \textit{and} corresponds to a genuine subword boundary in NLLB's vocabulary. Boundaries that span a single NLLB vocabulary item are removed. This reduces double-tokenization while retaining meaningful morphological signal.

\subsection{Morfessor-BPE}

Following \citet{banerjee2018morpheme}, we apply BPE \textit{within} Morfessor morphemes. First, Morfessor segments the corpus; then a BPE model (vocabulary size 4,000) is trained on the resulting morpheme tokens and applied within morpheme boundaries. This combines morphological segmentation with BPE's data-driven compression.

\subsection{UnigramLM}

We train a SentencePiece Unigram language model \cite{kudo2018subword} on the Mapudungun training data using the same 95\%-character-coverage vocabulary size heuristic as Mono BPE. Unlike BPE, UnigramLM directly optimizes a probabilistic segmentation model and can produce multiple segmentations for the same word, which has been shown to improve robustness in low-resource settings \cite{kudo2018subword}.

\subsection{Fertility Comparison}

Table~\ref{tab:fertility} shows the fertility (pre-segmented tokens per whitespace-delimited word) for each condition on the Mapudungun training side. Morfessor and Morfessor-VC have the lowest fertility (1.13), close to one-to-one morpheme--word mapping. BPE-based methods are more aggressive (1.55--2.04), and UnigramLM (1.98) is comparable to Mono BPE. The near-identical fertility of Morfessor and Morfessor-VC confirms that the VC merging step removes boundaries without adding new ones.

\begin{table}[h]
\centering
\begin{tabular}{lc}
\hline
\textbf{Condition} & \textbf{Fertility} \\
\hline
Standard BPE (NLLB) & --- \\
Joint BPE           & 1.66 \\
Mono BPE            & 2.04 \\
Optuna BPE          & 1.55 \\
Morfessor           & 1.13 \\
Morfessor-VC        & 1.13 \\
Morfessor-BPE       & 1.63 \\
UnigramLM           & 1.98 \\
\hline
\end{tabular}
\caption{Pre-segmentation fertility (tokens per whitespace word) on the Mapudungun training set. Standard BPE fertility is computed internally by NLLB and not directly comparable.}
\label{tab:fertility}
\end{table}

\section{Experimental Setup}
% TODO: Models: NLLB-200-distilled-600M, 1.3B, 3.3B
% TODO: Fine-tuning: 10 epochs, early stopping patience 3, lr=3e-5, batch 128, A100
% TODO: Evaluation: chrF++ (primary), BLEU (secondary), human evaluation (rank preference, Standard BPE vs.\ Morfessor-VC, 50 sentences × 2 directions)

\todo{Draft Experimental Setup}

\section{Results}
% TODO: Main results table: all conditions x 3 sizes x 2 directions (chrF++)
% TODO: Highlights: all morphology-aware methods cluster closely (43.1–45.8 arn->es); Morfessor-VC has slight edge; Standard BPE lags by ~2-3 chrF++ at 600M, gap closes at 3.3B
% TODO: Scaling: 600M Morfessor-VC \approx 3.3B standard BPE (efficiency result)
% TODO: vs LLMs: fine-tuned 600M beats Aya Expanse 8B by ~\textbf{[PLACEHOLDER]} chrF++

\todo{Insert main results table and discuss highlights/scaling}

\subsection{Human Evaluation}
% TODO: Human evaluation results (adequacy + fluency rankings)

\todo{Insert human evaluation results once collected}

\section{Analysis}

\subsection{Why Morfessor-VC Works}
% TODO: vocabulary coverage rate (\% of morpheme boundaries removed by VC step); fertility comparison table

\todo{Draft Morfessor-VC analysis and insert coverage/fertility table}

\subsection{Linguistic Analysis}
% TODO: qualitative annotation of ~25 arn->es sentences — morphological wellformedness, code-switching intrusion in es->arn output, dialect/register defaults; error taxonomy with 2–3 illustrative examples

\todo{Draft linguistic analysis based on Rogers' annotations}

\subsection{Code-switching and es$\rightarrow$arn Quality}
\label{sec:cs}
% TODO: automated CS detection on reference vs. output; correlation with per-sentence chrF++

\todo{Draft CS analysis and correlation}

\section{Conclusion}
% TODO: Summarize findings: Morfessor-VC success, scaling efficiency, and establishing a new baseline for Mapudungun MT.

\todo{Draft Conclusion}

\section*{Limitations}
% TODO: Discuss limitations (e.g., domain specificity of the AVENUE corpus, potential biases, reliance on Spanish as a pivot/bridge language, etc.)

\todo{Draft Limitations}

\section*{Acknowledgments}
% TODO: Acknowledge funders, MATRIX Lab, compute resources, and human evaluators.

\todo{Draft Acknowledgments}

% Custom bibliography entries only
\bibliography{custom}

\end{document}